\section{Physics-Informed Neural Network for Thermoelastic Prediction}

\subsection{Motivation for a physics-informed surrogate}

Among the four predictive routes considered in this thesis, the
physics-informed neural network (PINN) is the only one whose training
objective explicitly contains residuals of the governing thermoelastic
equations introduced in Chapter~3. The other three routes (FNO,
MeshGraphNet and the Transformer-style operator) are data-driven
surrogates: they are trained against simulation snapshots, but they do
not directly penalise violations of the coupled heat--momentum system.
The PINN component therefore plays a dual role inside the comparison
framework: it is both a candidate predictor and the only baseline that
embeds physical knowledge inside the loss.

\subsection{Network mapping and inputs}

The PINN is a single multilayer perceptron $f_{\theta}$ that takes a
point query in space--time together with the local material parameters
and returns the predicted thermoelastic state at that point:
\begin{equation}
f_{\theta}\colon\; 
\bigl(x,\,y,\,t,\,E,\,\nu,\,\rho,\,\alpha,\,k,\,C_p\bigr)
\;\longmapsto\;
\bigl(T,\,u,\,v\bigr).
\end{equation}
Here $(x,y)$ is the in-plane probe coordinate in the 2D plane-strain
domain, $t$ is the observation time, and the six material parameters
follow Section~3.4. The output triple $(T,u,v)$ contains the
temperature and the two in-plane displacement components; stress and
strain are recovered analytically from the predicted displacement and
temperature during loss assembly.

The network is implemented as an MLP with $L$ hidden layers of width
$d_h$ and a smooth activation $\varphi$ (the implementation uses
$\tanh$). Layer $\ell$ computes
\begin{equation}
h^{(\ell)} = \varphi\bigl(W^{(\ell)} h^{(\ell-1)} + b^{(\ell)}\bigr),
\qquad
h^{(0)} = z = (x,y,t,E,\nu,\rho,\alpha,k,C_p),
\end{equation}
and the output head produces
$(T,u,v) = W^{(L+1)} h^{(L)} + b^{(L+1)}$.
The smooth activation matters: $\tanh$ has bounded derivatives of
every order, so the higher-order derivatives needed by the PDE loss
(below) remain numerically stable.

\subsection{Composite training loss}

The PINN parameters $\theta$ are obtained by minimising a weighted sum
of four loss components on a batch of collocation points
$\{(x_i, y_i, t_i)\}$ drawn from the training scenarios:
\begin{equation}
\mathcal{L}(\theta)
=
\lambda_{\mathrm{sup}}\,\mathcal{L}_{\mathrm{sup}}
+ \lambda_{\mathrm{vel}}\,\mathcal{L}_{\mathrm{vel}}
+ \lambda_{\mathrm{el}}\,\mathcal{L}_{\mathrm{elastic}}
+ \lambda_{\mathrm{th}}\,\mathcal{L}_{\mathrm{thermal}}.
\end{equation}
The implementation-level assembly of this composite objective is shown
in Appendix~\ref{app:pinn-loss}.

Supervised loss $\mathcal{L}_{\mathrm{sup}}$ is the standard
mean-squared error between the predicted $(T,u,v)$ and the
ground-truth values exported from COMSOL at the corresponding
$(x_i,y_i,t_i)$.

Velocity consistency $\mathcal{L}_{\mathrm{vel}}$ enforces
that the displacement time derivative obtained by automatic
differentiation matches the velocity field that COMSOL reported as
$u_t,v_t$ in the same probe samples. This term prevents the network
from drifting toward spatially correct but temporally inconsistent
displacement fields.

Elastic-wave residual $\mathcal{L}_{\mathrm{elastic}}$
evaluates the residual of the linear-momentum balance
\begin{equation}
\rho\,\partial_t^2 u
\;-\;
\nabla\!\cdot\!\sigma
\;=\;0,
\qquad
\sigma
= \lambda\,(\nabla\!\cdot\!u)\,I
+ 2\mu\,\varepsilon
- \gamma\,\theta\,I,
\end{equation}
at the same collocation points, with $\varepsilon$ computed from the
predicted displacement gradients, $\theta = T - T_{\mathrm{ref}}$ the
temperature perturbation, and $\gamma = 3K\alpha$ the thermoelastic
coupling coefficient. All derivatives entering this residual are
obtained via automatic differentiation of the network output with
respect to its inputs.

Coupled-thermal residual $\mathcal{L}_{\mathrm{thermal}}$
evaluates the residual of the heat equation extended by the
thermoelastic coupling source:
\begin{equation}
\rho\,C_p\,\partial_t T
\;-\;
k\,\nabla^2 T
\;+\;
\gamma\,T_0\,\partial_t\!\bigl(\nabla\!\cdot\!u\bigr)
\;=\;0.
\end{equation}
The last term is exactly the volumetric-strain feedback discussed in
Section~3.5. Including it makes the PINN sensitive to the geometric
coupling between mechanical and thermal evolution; data-driven
surrogates without this term cannot detect it.

\subsection{Training data and sampling}

The PINN consumes the same per-rock 2D structured datasets that
Chapter~5 exported from COMSOL ($230$ nodes per material,
$101$ time samples). Each training row carries the probe coordinates,
the time stamp, the material parameters of the current rock, and the
ground-truth $(T,u,v)$ snapshot. The four loss terms are evaluated on
the same batch: supervised terms on the actual exported samples,
residuals on a randomly drawn collocation subset.

\subsection{Role in the comparison framework}

Because the PINN penalises violations of the governing equations
directly, its predictions are constrained to remain on a physically
meaningful manifold even when the supervised signal is noisy. The
comparison results of Chapter~12 will show that this is the only route
whose travel-time predictions stay within $10\%$ of the analytical
reference $r/V_p$ for every scenario in the balanced sweep, and the
only route whose time-series traces show a smooth time-dependent
drift at the probe. The other three routes are reported alongside the
PINN under the same v2 input/output contract, but they should be
read as data-driven baselines rather than physics-informed solvers.
